\subsection{Modular specialization}
\label{subsec:modular-specialization}

Suppose in this subsection that $\cC$ is modular.  Let $g_\eta$ be the
invertible simple object corresponding to $\eta$ under
\cite[Lemma~8.22.9(iii)]{EGNO2015}; explicitly,
$\alpha(g_\eta)=\eta$, where $\alpha$ is the closed positive-monodromy
automorphism in \cite[Equation~(8.95)]{EGNO2015}.  Put
\[
  \tau_+(\cC):=\tau_\eta^+(0),
  \qquad
  \tau_-(\cC):=\tau_\eta^-(0).
\]

\begin{theorem}[Modular closed formula]
\label{prop:modular-formula}
For every closed, oriented, simply connected smooth four-manifold $X$ and
every $\alpha\in H_2(X;\mathbb Z)$,
\[
  \CY_{\cC,\eta}(X,\alpha)
  =\operatorname{Dim}(\cC)\,
   \tau_+(\cC)^{b_2^+(X)}
   \tau_-(\cC)^{b_2^-(X)}
   \theta_{g_\eta}^{-Q_X(\alpha,\alpha)}.
\]
\end{theorem}

\begin{proof}
For every $n\in\mathbb Z$, Lemma~8.22.9(ii) of \cite{EGNO2015} gives
$\alpha(g_\eta^n)=\eta^n$.  Closing the remaining strand in
Equation~(8.95) gives, for every simple $i$,
\[
  s_{i,g_\eta^n}=d_i d_{g_\eta^n}\eta_i^n.
\]
Substitute this identity into the positive and negative $S$-matrix Gauss
identities, Equations~(8.58) and~(8.61) of \cite{EGNO2015}, and cancel the
nonzero quantum dimension of the invertible object $g_\eta^n$.  This gives
\[
  \tau_\eta^+(n)=\tau_+(\cC)\theta_{g_\eta^n}^{-1},
  \qquad
  \tau_\eta^-(n)=\tau_-(\cC)\theta_{g_\eta^n}.
\]
By the balancing equation and ribbon duality
\cite[Definition~8.10.1]{EGNO2015}, the twist restricted to the Picard group
is the quadratic form described immediately before
\cite[Definition~8.4.1]{EGNO2015}.  Hence
$\theta_{g_\eta^n}=\theta_{g_\eta}^{n^2}$, and therefore
\[
  \tau_\eta^+(n)=\tau_+(\cC)\theta_{g_\eta}^{-n^2},
  \qquad
  \tau_\eta^-(n)=\tau_-(\cC)\theta_{g_\eta}^{n^2};
\]
For an odd form, apply \cref{prop:odd-formula}.  Multiplying these identities
over its diagonal basis gives the stated power of $\theta_{g_\eta}$ and the
exponents $b_2^\pm(X)$; the negative coordinates occur as $-b_s$ there, but
$(-b_s)^2=b_s^2$.

For an even form, blow up once: $X\#\mathbb{CP}^2$ has odd intersection
form.  Apply the odd calculation to the marking $\alpha\oplus0$ and compare
it with the connected-sum formula.  The additional factor is
$\tau_+(\cC)$, which is nonzero in a modular category
\cite[Proposition~8.15.4]{EGNO2015}; cancelling it yields the same formula
for $X$.
\end{proof}

\subsubsection{Examples}
\label{subsubsec:modular-examples}

The following examples record exact consequences of
\cref{prop:modular-formula}.
\par\medskip

\begin{example}[The projective planes]
\label{ex:projective-planes}
Let $e^2=1$ generate $H_2(\mathbb{CP}^2;\mathbb Z)$ and let
$\bar e^2=-1$ generate $H_2(\overline{\mathbb{CP}}^{\,2};\mathbb Z)$.  For
$n\in\mathbb Z$,
\[
  \widehat{\CY}_{\cC,\eta}(\mathbb{CP}^2,ne)
    =\tau_+(\cC)\theta_{g_\eta}^{-n^2},
  \qquad
  \widehat{\CY}_{\cC,\eta}(\overline{\mathbb{CP}}^{\,2},n\bar e)
    =\tau_-(\cC)\theta_{g_\eta}^{n^2}.
\]
Thus the entire marked function on either projective plane is determined by
one Gauss sum and the twist of $g_\eta$.
\end{example}

\begin{example}[Diagonal rational surfaces]
\label{ex:diagonal-rational-surfaces}
For
\[
 X=p\mathbb{CP}^2\#q\overline{\mathbb{CP}}^{\,2},
 \qquad
 \alpha=(a_1,\ldots,a_p;b_1,\ldots,b_q),
\]
the modular closed formula gives
\[
 \widehat{\CY}_{\cC,\eta}(X,\alpha)
 =\tau_+(\cC)^p\tau_-(\cC)^q
  \theta_{g_\eta}^{-
    (\sum_r a_r^2-\sum_s b_s^2)}.
\]
In particular, for $\mathbb{CP}^2\#2\overline{\mathbb{CP}}^{\,2}$ the
exponent is $-(a^2-b^2-c^2)$.
\end{example}

\begin{example}[The hyperbolic block]
\label{ex:s2xs2}
Write $H=S^2\times S^2$ and choose the hyperbolic basis $u,v$, with
$u^2=v^2=0$ and $u\cdot v=1$.  Then
\[
  \widehat{\CY}_{\cC,\eta}(H,au+bv)
  =\tau_+(\cC)\tau_-(\cC)\theta_{g_\eta}^{-2ab}.
\]
For the zero marking $a=b=0$, this is precisely the factor introduced by a
hyperbolic stabilization and cancelled in the stabilization argument.
\end{example}

\begin{example}[$K3$]
\label{ex:k3-modular}
With the complex orientation, $b_2^+(K3)=3$ and $b_2^-(K3)=19$; hence
\[
  \widehat{\CY}_{\cC,\eta}(K3,\alpha)
  =\tau_+(\cC)^3\tau_-(\cC)^{19}
   \theta_{g_\eta}^{-Q_{K3}(\alpha,\alpha)}.
\]
No Kirby-diagram contraction is needed once the intersection form and the
self-intersection of $\alpha$ are known.
\end{example}

Independent numerical checks of these formulas, together with the
runtime, cloud-cost, classical-scaling, and quantum-query estimates for the
direct evaluator, are collected in
Appendix~\ref{app:computational-record}.

